\documentclass[12pt,a4paper]{article}

\usepackage{amsmath, amssymb, graphicx}
\usepackage{geometry}
\usepackage{setspace}
\usepackage{hyperref}

\geometry{margin=1in}
\setstretch{1.2}

\title{\textbf{Lecture Scribe}\\Fundamentals of Probability in Computing}
\author{
Name: Sakshi Patel\\
Student ID: AU2440206\\
Course: Fundamentals of Probability in Computing\\
}
\begin{document}
\date{}
\maketitle
\newpage

\section*{Step 1: Main Topics Identified (from the Lecture)}

The lecture revolves around three tightly connected ideas:

\begin{enumerate}
\item Transformation of Random Variables  
How probability changes when we apply a function to a random variable.

\item Function of Two Random Variables  
When output depends on two uncertain quantities simultaneously.

\item Illustrative Derivation  
Detailed case:
\[
Z = X + Y
\]
\end{enumerate}

Core message: Random variables are often transformed, and probability must be tracked after transformation.

\section*{2) Transformation of Random Variables}

\subsection*{2.1 Concept (Beginner View)}

Suppose we know the behavior of a variable $X$, but we use something derived from it.

Example:

Data size $X$  

Processing time:
\[
Y = 2X + 5
\]

We need the distribution of $Y$.  
This process is called transformation.

\subsection*{2.2 Definition}

A transformation occurs when:
\[
Y = g(X)
\]

Where:

$X$ → original random variable  

$g(\cdot)$ → function  

$Y$ → transformed variable  

Goal: find distribution of $Y$.

\subsection*{2.3 Intuition Behind Transformation}

Probability is conserved.

If a value of $X$ occurs with a certain likelihood, the corresponding value of $Y$ carries the same probability.

\[
P(X \text{ in interval}) = P(Y \text{ in mapped interval})
\]

\subsection*{2.4 Mathematical Derivation}

Assume continuous random variable.

Step 1: Known pdf
\[
f_X(x)
\]

Step 2: Transformation
\[
Y = g(X)
\]

Step 3: Inverse function
\[
X = g^{-1}(Y)
\]

Step 4: Probability conservation
\[
P(y \le Y \le y + dy) = P(x \le X \le x + dx)
\]

Thus:
\[
f_Y(y) dy = f_X(x) dx
\]

Rearranging:
\[
f_Y(y) = f_X(x)\left|\frac{dx}{dy}\right|
\]

This is the transformation formula.

\subsection*{2.5 Concept Check}

\begin{itemize}
\item Mapping intervals is valid
\item Probability is preserved
\item Derivative converts scale
\end{itemize}

\subsection*{2.6 Example}

Let:
\[
Y = X^2
\]
Assume:
\[
X \sim \text{Uniform}(-1,1)
\]

Inverse:
\[
X = \pm \sqrt{Y}
\]

Derivative:
\[
\frac{dx}{dy} = \frac{1}{2\sqrt{Y}}
\]

Total pdf:
\[
f_Y(y) = f_X(\sqrt{y})\left|\frac{dx}{dy}\right|
+ f_X(-\sqrt{y})\left|\frac{dx}{dy}\right|
\]

Reason: two values of $X$ produce the same $Y$.

\section*{3) Function of Two Random Variables}

Now uncertainty comes from two variables:

\[
Z = g(X,Y)
\]

\subsection*{3.1 Intuition}

Real systems depend on multiple uncertain factors:

\begin{itemize}
\item Network delay = transmission + queuing
\item ML error = noise + data bias
\item Battery life = temperature + usage
\end{itemize}

\subsection*{3.2 Definition}

\[
Z = g(X,Y)
\]

Goal: find distribution of $Z$.

\subsection*{3.3 Joint Distribution}

Start with:
\[
f_{X,Y}(x,y)
\]

This represents probability that:
\[
X=x \text{ and } Y=y
\]

\subsection*{3.4 Transformation Logic}

Coordinate transformation:
\[
(X,Y) \rightarrow Z
\]

Probability mass remains conserved.

\section*{4) Important Example: $Z = X + Y$}

\subsection*{4.1 Problem}

Given:
\[
Z = X + Y
\]

Find pdf of $Z$.

\subsection*{4.2 Conceptual Understanding}

Fix value $z$.

Find all combinations:
\[
x + y = z
\]

This forms a straight line in the $XY$ plane.

Every point contributes probability.

\subsection*{4.3 Mathematical Derivation}

\[
f_Z(z) = \int f_{X,Y}(x, z-x)\,dx
\]

If independent:
\[
f_{X,Y}(x,y) = f_X(x)f_Y(y)
\]

Then:
\[
f_Z(z) = \int f_X(x)f_Y(z-x)\,dx
\]

This operation is called convolution.

\subsection*{4.4 Why Convolution Works}

\begin{itemize}
\item Each pair $(x,y)$ producing $z$ contributes probability
\item Add all contributions → integration
\item Independence → multiplication of pdfs
\end{itemize}

\subsection*{4.5 Example}

Let:
\[
X,Y \sim \text{Uniform}(0,1)
\]

Then:
\[
f_Z(z)=
\begin{cases}
z, & 0<z<1 \\
2-z, & 1<z<2
\end{cases}
\]

Shape: triangular distribution.

\section*{5) Conceptual + Mathematical + Intuitive Summary}

\begin{center}
\begin{tabular}{|c|p{8cm}|}
\hline
Level & Understanding \\
\hline
Conceptual & Probability stays conserved after transformation \\
Mathematical & Use derivative and integration \\
Intuitive & Map original values to new values \\
Computational & Used in ML, networks, simulations \\
\hline
\end{tabular}
\end{center}

\section*{6) Complex Ideas Explained}

\begin{itemize}
\item Why derivative appears? → Scale changes during mapping.
\item Why inverse needed? → Express pdf in new variable.
\item Why integration? → Many inputs produce same output.
\item Why independence matters? → Simplifies joint distribution.
\end{itemize}

\section*{7) Concise Lecture Summary}

\begin{itemize}
\item Random variables are transformed using functions.
\item Distribution must be recomputed.
\item Single variable → derivative rule.
\item Two variables → joint distribution.
\item Sum → convolution.
\end{itemize}

\section*{8) Key Formula Sheet}

Transformation:
\[
f_Y(y) = f_X(x)\left|\frac{dx}{dy}\right|
\]

Joint pdf:
\[
f_{X,Y}(x,y)
\]

Sum of RVs:
\[
f_Z(z)=\int f_X(x)f_Y(z-x)\,dx
\]

\section*{9) Probable Exam Questions}

Conceptual:
\begin{itemize}
\item Explain transformation of random variables.
\item Why derivative is used in pdf transformation?
\end{itemize}

Derivations:
\begin{itemize}
\item Derive pdf of $Y=g(X)$.
\item Derive pdf of $Z=X+Y$.
\end{itemize}

Numerical:
\begin{itemize}
\item Uniform + Uniform
\item Normal + Normal
\item Square transformation
\end{itemize}

Theory:
\begin{itemize}
\item Joint distribution
\item Independence
\end{itemize}

\section*{10) Key Insight}

Students often memorize formulas but struggle with derivations.

Examiners typically ask:

\begin{itemize}
\item Why is each step valid?
\item Where did the derivative come from?
\item Why is integration required?
\end{itemize}

Understanding the logic behind transformations leads to better performance.

\end{document}
